\GalSourcePageBoundary{077}{L0076}{48}{48}

On raisonnera sur ce groupe comme sur le précédent, et il s'ensuivra
que le premier groupe dans l'ordre des décompositions,
c'est-à-dire le groupe \emph{actuel} de l'équation, ne peut contenir que
des substitutions de la forme
\[
x_{k},\qquad x_{ak+b}.
\]

Donc, {\itshape si une équation irréductible de degré premier est
soluble par radicaux, le groupe de cette équation ne saurait
contenir que des substitutions de la forme
\[
x_{k},\qquad x_{ak+b},
\]
$a$ et $b$ étant des constantes.}

Réciproquement, si cette condition a lieu, je dis que l'équation
sera soluble par radicaux. Considérons, en effet, les fonctions
\[
\begin{aligned}
&(x_{0}+\alpha x_{1}+\alpha^{2}x_{2}+\ldots+\alpha^{n-1}x_{n-1})^{n}=X_{1},\\
&(x_{0}+\alpha x_{a}+\alpha^{2}x_{2a}+\ldots+\alpha^{n-1}x_{(n-1)a})^{n}=X_{a},\\
&(x_{0}+\alpha x_{a^{2}}+\alpha^{2}x_{2a^{2}}+\ldots+\alpha^{n-1}x_{(n-1)a^{2}})^{n}=X_{a^{2}},\\
&\hspace{12em}\ldots
\end{aligned}
\]
$\alpha$ étant une racine $n^{\text{ième}}$ de l'unité, $a$ une racine primitive de $n$.

Il est clair que toute fonction invariable par les substitutions
circulaires des quantités $X_{1},X_{a},X_{a^{2}},\ldots$ sera, dans ce cas, immédiatement
connue. Donc, on pourra trouver $X_{1},X_{a},X_{a^{2}},\ldots$, par
la méthode de M. Gauss pour les équations binômes. Donc, etc.

Ainsi, pour qu'une équation irréductible de degré premier soit
soluble par radicaux, il \emph{faut} et il \emph{suffit} que toute fonction invariable
par les substitutions
\[
x_{k},\qquad x_{ak+b}
\]
soit rationnellement connue.

Ainsi, la fonction
\[
(X_{1}-X)(X_{a}-X)(X_{a^{2}}-X)\ldots
\]
devra, quel que soit $X$, être connue.
